% ==============================================================================
% ПАРТИЯ B03-016: DISS-005-B0698 .. DISS-005-B0734
% Глава 1: Поведение органов БАВР при феррорезонансных процессах при дуговых перемежающихся ОЗЗ
% Замечания: DISS-005-C0070 .. DISS-005-C0073
% Иллюстрации: image48.png .. image53.png
% ==============================================================================

% DISS-005-B0698
\subsection{Поведение органов БАВР при феррорезонансных процессах при дуговых перемежающихся \commentnote{DISS-005-C0070}{Евдаков Алексей}{Тут по большей части разбор возникновения блокировки по БНН. К РНМ можно притянуть в том контексте, что он тоже использует напряжение.\newline\newlineНо также, очень актуально будет упомянуть и описать в контексте статистки по ОЗЗ, что возникает перенапряжение, но оно не реальное.}{\commentnote{DISS-005-C0071}{Алексей Евдаков [2]}{И надо будет обязательно перечитать, так как пока что, это прост скопированный текст.}{\commentnote{DISS-005-C0072}{Евдаков Алексей}{Данные файлы находятся в папке\newline«\_вспомогательные файлы\textbackslash{}Неман»\newline}{ОЗЗ}}}}
\label{sec:ch1_ferroresonance_ozz}

% DISS-005-B0699
\subsubsection{Информация о применении и актуальности проработки темы}

% DISS-005-B0700
Приводится разбор аварийного процесса на ТП-83 «Неман» от 15.10.2025 в 15:03. Установлен терминал БАВР производства ООО «АПС».

% DISS-005-B0701
В сети произошло однофазное замыкание на землю на фазе В (до 2,5~с на осциллограмме), которое позже переросло в двухфазное замыкание (АВ, длилось 0,57~с), что привело к аварийному снижения напряжения. На рисунке 1 представлено общее представление протекания аварии на 2 секции шин.

% DISS-005-B0702
Можно бы здесь или в другом месте добавить вот эту статью, в ней описано то, как много случается ОЗЗ, по каким причинам и т.д. А в рамках этого учатска, или где-то в другом месте (выводах) добавить о том, что часто БАВР срабатывание как раз из-за ОЗЗ, которые переходят в двухфазное

% DISS-005-B0703
Вот статья, есть у меня в списке

% DISS-005-B0704
(\commentnote{DISS-005-C0073}{Алексей Евдаков [2]}{analiz\_statistiki\_tehnologicheskih\_narusheniy\_v\_kabelnyh\_raspredelitelnyh\newline}{Анализ }статистики технологических нарушений в кабельных сетях 6(10)кВ) .

% DISS-005-B0706 .. DISS-005-B0707 (DISS-005-FIG-045 / DISS-005-TAB-031)
\begin{figure}[htbp]
    \centering
    \includegraphics[width=0.85\textwidth]{image48.png}
    \caption{Рисунок 1. Развитие аварии на второй секции шин.}
    \label{fig:ch1_img48}
\end{figure}

% DISS-005-B0708
При данной просадке напряжения должно было произойти срабатывание устройства, чего не случилось ввиду возникновения ранее запрета по БНН (блокировка при неисправностях в цепях напряжения). Возникновение данного сигнала можем видеть на рисунке 2. Это привело к блокировке работы всех защит, связанных с напряжением (27/27, 59, РНМ, пуск по углу) на временной промежуток с 1,3~с до 4,1~с, а следовательно и пропаже готовности БАВР.

% DISS-005-B0709 .. DISS-005-B0710 (DISS-005-FIG-046 / DISS-005-TAB-032)
\begin{figure}[htbp]
    \centering
    \includegraphics[width=0.85\textwidth]{image49.png}
    \caption{Рисунок 2. Запрет работы БАВР по сигналу БНН.}
    \label{fig:ch1_img49}
\end{figure}

% DISS-005-B0712
Причиной возникновения данной блокировки скорее всего* послужила активации функции БНН по обратной последовательности напряжения ($U_2$), которая предназначена для регистрации обрыва во вторичных цепях напряжения.

% DISS-005-B0713
*Примечание: Исходного файла конфигурации не имеем, поэтому утверждать не будем, но вероятность крайне высока.

% DISS-005-B0714
Подобная несиметрия сигнала обусловлена процессами, протекающими во внешней цепи по отношению к терминалу. Скорее всего в ТН, находящимся на секции шин, возникли резонансные процессы от протекания дуговых перемежающихся ОЗЗ, которые привели к радикальному искажению сигнала на фазе С (рисунок 3), по фазам А было 110~В (стандартное линейное), по фазам В было 26~В (как раз замыкание фазы на землю), на фазе С на ТН СШ показывается значение в 170~В (выходящее за ожидаемые диапазоны), а на ТН КЛ 84~В (как раз ожидаемые значения).

% DISS-005-B0715 .. DISS-005-B0716 (DISS-005-FIG-047 / DISS-005-TAB-033)
\begin{figure}[htbp]
    \centering
    \includegraphics[width=0.85\textwidth]{image50.png}
    \caption{Рисунок 3. Фазные напряжения.}
    \label{fig:ch1_img50}
\end{figure}

% DISS-005-B0717
Более наглядно это видно при оценке линейных значений (рисунок 4) -- при оценке измерений с ТН КЛ, линейные напряжения стабильны и составляют 104~В, а как раз на ТН СШ -- произошёл перекос ($U_{ab}=103$~В, $U_{ab}=173$~В, $U_{ab}=112$~В). И как раз это привело к возникновению измеренной обратной последовательности (28~В) и сигналу БНН.

% DISS-005-B0718 .. DISS-005-B0719 (DISS-005-FIG-048 / DISS-005-TAB-034)
\begin{figure}[htbp]
    \centering
    \includegraphics[width=0.85\textwidth]{image51.png}
    \caption{Рисунок 4. Линейные напряжения и напряжения обратной последовательности.}
    \label{fig:ch1_img51}
\end{figure}

% DISS-005-B0721
Отметим, что после завершения аварийных процессов напряжения на ТН СШ вернулись в норму и по всем фазам значения были около 60~В (рисунок 5).

% DISS-005-B0722 .. DISS-005-B0723 (DISS-005-FIG-049 / DISS-005-TAB-035)
\begin{figure}[htbp]
    \centering
    \includegraphics[width=0.85\textwidth]{image52.png}
    \caption{Рисунок 5. Послеаварийное измерение.}
    \label{fig:ch1_img52}
\end{figure}

% DISS-005-B0725
На рисунке 6 более близко представлено протекание дугового перемежающегося ОЗЗ. Можем видеть возникновение пробоев в случайные моменты времени, что приводит различным колебательным процессам в сети.

% DISS-005-B0726 .. DISS-005-B0727 (DISS-005-FIG-050 / DISS-005-TAB-036)
\begin{figure}[htbp]
    \centering
    \includegraphics[width=0.85\textwidth]{image53.png}
    \caption{Рисунок 6. Протекание дугового перемежающегося замыкания на землю.}
    \label{fig:ch1_img53}
\end{figure}

% DISS-005-B0729 .. DISS-005-B0732
Выводы:
\begin{itemize}[leftmargin=*,itemsep=0pt]
    \item Причиной блокировки работы БАВР послужило некорректное значение в измерительных цепях с ТН СШ по фазе С ввиду резонансных явления вызванных дуговым перемежающимся ОЗЗ.
    \item Для исключения данной блокировки предлагается отключить БНН по $U_2$.
    \item Так же рекомендуется в будущем использовать ТН с антирезонансными свойствами и/или разобраться с причинами возникновения резонансных процессов в текущей сети.
\end{itemize}
